\documentclass[12pt, a4paper]{article}
\usepackage[utf8]{inputenc}
\usepackage{ctex} 
\usepackage{geometry}
\usepackage{fancyhdr}
\usepackage{setspace}
\usepackage{tabularx} 
\usepackage{array}

% 全局页面设置
\geometry{top=3cm, bottom=3cm, left=2.5cm, right=2.5cm}
\pagestyle{empty} 

% 定义普通正文方框命令
\newcommand{\mylabsection}[2]{%
    \par\nointerlineskip%
    \noindent%
    \begin{tabularx}{\textwidth}{|@{}X@{}|}% 使用 @{} 消除列间距，确保宽度一致
        \hline
        \hspace{1em}{\zihao{5} #1} \\
        \rule{0pt}{#2} \\ \hline
    \end{tabularx}%
}

% 定义全页完整方框命令（仅此页左下角带标识）
\newcommand{\myfullpagebox}{%
    \newpage%
    \noindent%
    \begin{tabularx}{\textwidth}{|@{}X@{}|} % 确保与普通框等宽
        \hline
        \rule{0pt}{22.5cm} \\ % 方框高度
        \hline
    \end{tabularx}%
    \par\nointerlineskip%
    \noindent{\zihao{5}\songti 深圳大学学生实验报告用纸}% 左下角对齐
}

\begin{document}

% --- 第一页：封面 (左边距增大) ---
\newgeometry{top=3cm, bottom=3cm, left=4cm, right=2.5cm} 

\begin{center}
    \vspace*{1cm}
    {\Huge \textbf{深 圳 大 学 实 验 报 告}} \\
    \vspace{2cm}
\end{center}

\begin{flushleft}
    \zihao{4} \bfseries 
    \begin{spacing}{1.2}
        \begin{tabular}{ll}
            课程名称：\underline{\makebox[25em][l]{}} \\[8ex]
            实验项目名称：\underline{\makebox[23em][l]{}} \\[8ex]
            学院：\underline{\makebox[25em][l]{}} \\[8ex]
            专业：\underline{\makebox[25em][l]{}} \\[8ex]
            指导教师：\underline{\makebox[25em][l]{}} \\[8ex]
            报告人：\underline{\makebox[6em][l]{}}\hspace{1.5em}学号：\underline{\makebox[6em][l]{}}\hspace{1.5em}班级：\underline{\makebox[6em][l]{}} \\[8ex]
            实验时间：\underline{\makebox[25em][l]{}} \\[8ex]
            实验报告提交时间：\underline{\makebox[23em][l]{}} 
        \end{tabular}
    \end{spacing}
\end{flushleft}

\vfill
\begin{center}
    {\Large \textbf{教务部制}}
\end{center}

\newpage 
\restoregeometry % 恢复全局边距

% --- 第二页：实验正文 ---
\mylabsection{实验目的与要求：}{10cm}%
\mylabsection{方法、步骤：}{10cm}%
\mylabsection{实验过程及内容：}{10cm}%

% 数据处理分析 
\mylabsection{数据处理分析：}{10cm}%

% --- 第三页：全页大方框 (仅在此页左下角有标识) ---
\myfullpagebox

% --- 第四页：结论与评分 ---
\mylabsection{实验结论：}{7cm}%

% 评分区
\par\nointerlineskip
\noindent
\begin{tabularx}{\textwidth}{|@{}X@{}|}
    \hline
    \hspace{1em}\zihao{-4} 指导教师批阅意见：\\[4cm] 
    \hspace{1em}成绩评定： \\[4cm] 
    \hspace*{23em} 指导教师签字： \\
    \hspace*{23em} \makebox[3em]{}年\makebox[2em]{}月\makebox[2em]{}日 \\ [1em]
\end{tabularx}%

% 备注框
\mylabsection{备注：}{2cm}%

% 底部注释
\begin{spacing}{1.2}
    \zihao{5}
    \noindent \begin{tabular}{@{}l p{35em}@{}}
        注： & 1、报告内的项目或内容设置，可根据实际情况加以调整和补充。 \\
             & 2、教师批改学生实验报告时间应在学生提交实验报告时间后 10 日内。
    \end{tabular}
\end{spacing}

\end{document}